\chapter[Fluid Mechanics]{Fluid Mechanics: Balance Laws, Navier-Stokes, And Stability}
\label{chap:fluid}

\section{What Is A Fluid?}
\label{sec:what-is-fluid}

A simple Newtonian fluid is a continuum that has lost its memory of a reference
shear. Push a finger through honey: the displaced material moves, stretches,
and eventually fills in behind the finger, but no permanent shear strain is
stored. Pull the finger out and the honey does not snap back. This is a
constitutive idealization, not the definition of every material called a fluid:
viscoelastic fluids, Bingham fluids, granular media, and superfluids require
different constitutive structures.

This is why the Eulerian description is convenient for many fluid problems. We
observe fields at fixed points in space: velocity, density, pressure, and
stress. The moving-particle viewpoint returns through the material derivative
and transport theorems, and remains essential for ideal-fluid variational
principles, vortex-line transport, mixing, and free-surface motion. The
historical names are literal here: Euler wrote the inviscid fluid equations in
Eulerian form in the 1750s, while the Lagrangian description follows labelled
material particles. The preceding elastic and deforming-body chapters kept the
material reference description in the foreground; this chapter shifts the
emphasis to Eulerian fields without treating them as the only natural
description.

The basic fields are functions on physical space:
\[
  \rho(x,t),\qquad u(x,t),\qquad p(x,t),
\]
where \(\rho\) is mass density, \(u\) is velocity, and \(p\) is pressure. A
fluid particle can still be labelled by a material coordinate \(X\), with
\[
  x=\varphi(X,t),\qquad u(\varphi(X,t),t)=\partial_t\varphi(X,t).
\]
The Eulerian view keeps \(x\) fixed and watches fluid pass through; the
Lagrangian view follows the particle labelled by \(X\).

\begin{assumption}[Continuum fluid]
The molecular scale is much smaller than the length scales being modelled, so
density, velocity, pressure, and stress are treated as smooth fields. Across
shocks, contact discontinuities, or interfaces, the differential equations must
be replaced by weak forms and jump conditions.
\end{assumption}

\paragraph{Roadmap.}
The derivation is staged in the order of the modelling choices. First we learn how to differentiate
quantities attached to moving fluid parcels. Then we convert integral balance
laws into local partial differential equations for mass and momentum. Only
after those kinematic and balance statements are in place do we introduce the
Newtonian constitutive law that turns Cauchy's momentum equation into the
Navier-Stokes equation. The later sections use that equation in three ways:
exact laminar solutions, stability analysis, and Hamiltonian point-vortex
reduction.

Elastic bodies use a material map as the primary unknown; deforming bodies add
frame choices and connections to material shape. A simple fluid still permits
material labels in derivations, but its closed equations are most useful as
equations for Eulerian fields: density, velocity, pressure, and stress at
points in space. The common continuum thread is stress and balance; the
constitutive question shifts from stored deformation energy to rate-dependent
stress.

The staging matters because the Navier-Stokes equation has three different
ingredients. Its inertial term is kinematics, its pressure and stress terms
come from momentum balance, and its viscous term comes from the Newtonian
constitutive law. Separating these ingredients makes it clear which part of the
model changes when one studies inviscid flow, non-Newtonian fluids, elasticity,
or low-Reynolds-number swimming.

\paragraph{Derivation audit trail.}
Every continuum equation below is paired with its conservation statement. The
material derivative comes from following one labelled parcel. The Reynolds
transport theorem converts derivatives of moving integrals into Eulerian fields.
Mass conservation comes from conserving mass in every material volume. Cauchy's
equation comes from momentum balance and the existence of a traction stress
tensor. Navier-Stokes appears after the Newtonian stress law is imposed, so
\(\rho Du/Dt=-\nabla p+\mu\Delta u\) records a sequence of modelling choices:
kinematics, mass balance, momentum balance, and constitutive law.

\section{Material Derivative}
\label{sec:material-derivative}

If \(f(x,t)\) is a scalar or vector field, its value along a particle path
\(x(t)=\varphi(X,t)\) changes as
\[
  \frac{\dd}{\dd t}f(x(t),t)
  =
  \partial_t f(x(t),t)+\dot x(t)\cdot\nabla f(x(t),t).
\]
Since \(\dot x=u(x(t),t)\), the material derivative is
\[
  \boxed{
  \frac{D f}{D t}
  =
  \partial_t f+u\cdot\nabla f.
  }
\]
For the velocity field itself,
\[
  \frac{D u}{D t}
  =
  \partial_tu+(u\cdot\nabla)u
\]
is the particle acceleration. The nonlinear term is kinematic: it appears even
before any force law is chosen, because the velocity field is sampled along a
moving trajectory.

\section{Reynolds Transport Theorem}
\label{sec:reynolds-transport}

Let \(V(t)\) be a material volume: it moves with the fluid, so its boundary
velocity is \(u\). For a scalar field \(f\), we want the time derivative of
\[
  I(t)=\int_{V(t)} f(x,t)\,\dd^3x.
\]
There are two sources of change. The value of \(f\) can change along each
particle path, and the volume element occupied by neighboring particles can
expand or contract. The Reynolds transport theorem is exactly the formula that
keeps both effects.
Reynolds gave the theorem in this systematic continuum-mechanics form in his
1903 treatise, but the pullback idea is already present in Euler's Lagrangian
description of fluids.

\begin{figure}[htbp]
\centering
\begin{tikzpicture}[scale=1.0, >=Latex, font=\small]
  \draw[->] (-0.4,0) -- (6.8,0) node[right] {\(x^1\)};
  \draw[->] (0,-0.4) -- (0,4.0) node[above] {\(x^2\)};
  \draw[blue!70!black, thick, fill=blue!5]
    (1.0,0.9) .. controls (1.8,0.35) and (3.2,0.55) .. (3.85,1.15)
    .. controls (4.6,1.85) and (4.15,3.0) .. (3.0,3.2)
    .. controls (1.75,3.45) and (0.75,2.45) .. (1.0,0.9);
  \draw[red!75!black, thick, ->] (1.3,1.0) -- (1.75,0.55);
  \draw[red!75!black, thick, ->] (3.9,1.45) -- (4.5,1.55);
  \draw[red!75!black, thick, ->] (2.95,3.15) -- (3.15,3.75);
  \draw[red!75!black, thick, ->] (0.9,2.4) -- (0.45,2.75);
  \node[blue!70!black] at (2.45,2.0) {\(V(t)\)};
  \node[red!75!black, fill=white, inner sep=2pt] at (5.28,3.10)
    {boundary velocity \(u\)};
  \draw[red!75!black, ->] (4.75,2.88) -- (3.15,3.50);
  \node[align=center, fill=white, inner sep=2pt] at (5.10,-0.32)
    {material points stay\\inside the moving volume};
  \draw[black!55, ->] (4.55,-0.05) -- (3.55,1.20);
\end{tikzpicture}
\caption{A material volume is carried by the fluid velocity. Its boundary
velocity is \(u\), so the time derivative of an integral over \(V(t)\) must
include both the material change of the integrand and the volume change
measured by \(\nabla\cdot u\).}
\label{fig:reynolds-material-volume}
\end{figure}

\begin{derivation}[Transport through a moving volume]
Write \(V(t)=\varphi(U,t)\) for a fixed material region \(U\). Let
\[
  F^i{}_A=\frac{\partial\varphi^i}{\partial X^A},
  \qquad
  J=\det F.
\]
Changing variables gives
\[
  I(t)=\int_U f(\varphi(X,t),t)J(X,t)\,\dd^3X.
\]
Differentiate:
\[
  \dot I
  =
  \int_U
  \left(
    \frac{D f}{D t}J+f\dot J
  \right)\dd^3X.
\]
The Jacobian identity is
\[
  \dot J=J\nabla\cdot u.
\]
For any smooth path \(F(t)\in\operatorname{GL}(n)\), Jacobi's formula is
\[
  \frac{\dd}{\dd t}\det F
  =
  \det F\,\tr(F^{-1}\dot F).
\]
Indeed,
\[
  \det(F+\delta F)
  =
  \det F\,\det(I+F^{-1}\delta F)
  =
  \det F\left(1+\tr(F^{-1}\delta F)+O(\delta F^2)\right).
\]
Applying this formula to \(J=\det F\) gives
\[
  \dot J=J\,\tr(F^{-1}\dot F),
\]
and \(\dot F^i{}_A=\partial_Au^i(\varphi,t)\), so
\(\tr(F^{-1}\dot F)=\partial_i u^i\). Returning to spatial variables,
\[
  \boxed{
  \frac{\dd}{\dd t}\int_{V(t)} f\,\dd^3x
  =
  \int_{V(t)}
  \left(
    \frac{D f}{D t}+f\nabla\cdot u
  \right)\dd^3x
  =
  \int_{V(t)}
  \left(
    \partial_t f+\nabla\cdot(fu)
  \right)\dd^3x.
  }
\]
\end{derivation}

The theorem separates two effects: the field \(f\) changes along particles, and
the material volume expands or contracts.

\section{Conservation Of Mass}
\label{sec:mass-conservation}

Mass in a material volume is constant:
\[
  \frac{\dd}{\dd t}\int_{V(t)}\rho\,\dd^3x=0.
\]
Using the transport theorem,
\[
  \int_{V(t)}
  \left(
    \partial_t\rho+\nabla\cdot(\rho u)
  \right)\dd^3x=0.
\]
Because this holds for every material volume,
\[
  \boxed{
  \partial_t\rho+\nabla\cdot(\rho u)=0.
  }
\]
Equivalently,
\[
  \frac{D\rho}{Dt}+\rho\nabla\cdot u=0.
\]
For constant-density incompressible flow,
\[
  \boxed{\nabla\cdot u=0.}
\]

Two statements are being separated here. Incompressibility means that material
volumes are preserved, so \(\nabla\cdot u=0\). Constant density is then
consistent with mass conservation, but it is not a separate equation of
motion. This distinction matters in stratified or variable-density flows, where
the velocity may be divergence-free in an approximation while density remains
an advected scalar.

\section{Cauchy Momentum Balance}
\label{sec:cauchy-momentum}

Cauchy's 1822--1823 stress principle is the common ancestor of the elastic and
fluid balance laws used in these notes. The same traction theorem that produced
stress in Chapter~\ref{chap:elastic-body} now produces the stress tensor of a
moving fluid.

Newton's second law for a material volume \(V(t)\) is
\[
  \frac{\dd}{\dd t}\int_{V(t)}\rho u\,\dd^3x
  =
  \int_{V(t)}\rho b\,\dd^3x
  +
  \int_{\partial V(t)} t(n)\,\dd S,
\]
where \(b\) is body force per unit mass and \(t(n)\) is surface traction on a
surface whose outward unit normal is \(n\).

\begin{proposition}[Cauchy's traction theorem]
Under the same locality and continuity hypotheses used in the tetrahedron
argument of Section~\ref{sec:spatial-stress-frame-indifference}, there is a
tensor \(\sigma(x,t)\) such that
\[
  t(n)=\sigma n,
\]
where \(\sigma\) is the Cauchy stress tensor.
\end{proposition}

Balance of angular momentum gives the symmetry of this tensor in the absence of
couple stresses. For a material volume \(V(t)\), moment balance about the
origin is
\[
  \frac{\dd}{\dd t}\int_{V(t)} \rho x\times u\,\dd^3x
  =
  \int_{V(t)} \rho x\times b\,\dd^3x
  +
  \int_{\partial V(t)} x\times(\sigma n)\,\dd S.
\]
Using the divergence theorem on the surface term gives the volume integral of
\(x\times(\nabla\cdot\sigma)\) plus the axial vector of
\(\sigma-\sigma^T\). Comparing with \(x\times\) the linear momentum balance
leaves only this axial term. Since the volume is arbitrary,
\[
  \sigma=\sigma^T.
\]
In indices, the extra term is
\[
  \epsilon_{ijk}\sigma_{kj}\,e_i,
\]
which is twice the axial vector of the skew tensor
\((\sigma-\sigma^T)/2\). Angular-momentum balance without internal body
couples forces this vector to vanish pointwise.

\begin{derivation}[Local momentum equation]
Apply the transport theorem to \(f=\rho u\). Using mass conservation, the
left-hand side becomes
\[
\begin{aligned}
  \frac{\dd}{\dd t}\int_{V(t)}\rho u\,\dd^3x
  &=
  \int_{V(t)}
  \left(
    \partial_t(\rho u)+\nabla\cdot(\rho u\otimes u)
  \right)\dd^3x\\
  &=
  \int_{V(t)}
  \rho\left(\partial_tu+u\cdot\nabla u\right)\dd^3x.
\end{aligned}
\]
The traction term is converted by the divergence theorem:
\[
  \int_{\partial V(t)}\sigma n\,\dd S
  =
  \int_{V(t)}\nabla\cdot\sigma\,\dd^3x.
\]
Therefore
\[
  \int_{V(t)}
  \left[
    \rho\frac{D u}{D t}
    -
    \nabla\cdot\sigma
    -
    \rho b
  \right]\dd^3x=0.
\]
Since \(V(t)\) is arbitrary,
\[
  \boxed{
  \rho\frac{D u}{D t}=\nabla\cdot\sigma+\rho b.
  }
\]
\end{derivation}

The localization step is the continuum analogue of the fundamental lemma of
the calculus of variations: if the integral of a smooth density vanishes for
every sufficiently small material volume, then the density itself must vanish
pointwise. This equation is not yet a closed fluid equation. It is a balance
law. Closure comes from a constitutive relation for \(\sigma\).

The same balance law will reappear in the elasticity chapter with a different
constitutive relation and a different choice of variables. That repetition is
intentional. Momentum balance is universal; the material model enters when one
decides how stress depends on deformation, rate of deformation, pressure,
temperature, or internal variables.

\section{Newtonian Stress And The Navier-Stokes Equation}
\label{sec:navier-stokes-derivation}

The name Navier-Stokes compresses a longer history: Navier introduced a
viscous molecular model in 1822, Cauchy supplied the stress framework,
Saint-Venant developed continuum viscous stresses, and Stokes gave the
standard nineteenth-century formulation in 1845. The modern axiomatic language
of objectivity and constitutive response was systematized by Truesdell and
Noll in the twentieth century.

The balance laws so far are universal for continua, but they are not yet a
closed fluid model: the stress tensor \(\sigma\) is still unknown. A
constitutive law supplies the missing relation between stress and motion. For a
simple Newtonian fluid, the stress may depend on the instantaneous rate of
deformation but not on the accumulated shear history. This is the point at
which fluids separate from elastic solids.

The stress of a simple Newtonian fluid is assumed to depend locally and
linearly on the instantaneous rate of deformation, not on the deformation
itself. The velocity gradient splits into symmetric and antisymmetric parts:
\[
  \nabla u=D+W,
  \qquad
  D_{ij}=\frac12(\partial_i u_j+\partial_j u_i),
  \qquad
  W_{ij}=\frac12(\partial_i u_j-\partial_j u_i).
\]
The antisymmetric part \(W\) represents local rigid rotation. It should not
produce viscous stress in an isotropic simple fluid. This statement is a
constitutive assumption with three ingredients:
\[
\begin{array}{ll}
\text{locality:} & \sigma(x,t) \text{ depends on } u \text{ through }
  \nabla u(x,t),\\
\text{objectivity:} & \text{superposed rigid rotations do not create stress},\\
\text{isotropy:} & \text{there is no preferred material direction}.
\end{array}
\]
Objectivity removes \(W\). Isotropy says that a linear map from the symmetric
tensor \(D\) to a symmetric stress can only use the deviatoric part of \(D\)
and its trace. Within the local, linear, objective, isotropic simple-fluid class
defined by these assumptions, the most general linear relation is therefore
\[
  \boxed{
  \sigma=-pI+2\mu D+\lambda(\nabla\cdot u)I.
  }
\]
Here \(p\) is pressure, \(\mu\geq0\) is shear viscosity, and \(\lambda\) is the
second viscosity coefficient. The symmetric strain-rate tensor decomposes as
\[
  D=D_0+\frac13(\tr D)I,\qquad \tr D_0=0.
\]
The representation statement used above is elementary in the linear case. A
linear isotropic map \(L\) on symmetric tensors must commute with every
rotation:
\[
  L(QDQ^T)=QL(D)Q^T.
\]
The one-dimensional subspace spanned by \(I\) is invariant, and the
five-dimensional space of traceless symmetric tensors is irreducible under
rotations. Hence \(L(D)=a(\tr D)I+bD_0\), which is the displayed law after
renaming constants.

For constant \(\mu,\lambda\),
\[
\begin{aligned}
  (\nabla\cdot\sigma)_i
  &=
  -\partial_i p
  +
  \mu\partial_j(\partial_i u_j+\partial_j u_i)
  +
  \lambda\partial_i(\partial_k u_k)\\
  &=
  -\partial_i p
  +
  \mu\Delta u_i
  +
  (\lambda+\mu)\partial_i(\nabla\cdot u).
\end{aligned}
\]
Thus the compressible Navier-Stokes momentum equation is
\[
  \boxed{
  \rho
  \left(
    \partial_tu+u\cdot\nabla u
  \right)
  =
  -\nabla p+\mu\Delta u+(\lambda+\mu)\nabla(\nabla\cdot u)+\rho b.
  }
\]
Together with mass conservation, an equation of state, and an energy or
temperature equation when needed, this describes a Newtonian compressible fluid.
The displayed equation combines three independent ingredients: Cauchy momentum
balance, the Newtonian stress law, and the constant-viscosity simplification. If
viscosity varies with temperature or composition, the divergence of
\(2\mu D+\lambda(\nabla\cdot u)I\) must be kept in product-rule form.

For incompressible flow with constant \(\rho\),
\[
  \nabla\cdot u=0,
\]
and pressure plays the role, in this incompressible model, of the Lagrange
multiplier enforcing this constraint. In compressible flow pressure is instead a
thermodynamic variable closed by an equation of state. The incompressible
equations reduce to
\[
  \boxed{
  \begin{aligned}
  \rho(\partial_tu+u\cdot\nabla u)&=-\nabla p+\mu\Delta u+\rho b,\\
  \nabla\cdot u&=0.
  \end{aligned}
  }
\]

Read the incompressible equations term by term. The material acceleration
\(\partial_tu+u\cdot\nabla u\) is the acceleration of a moving fluid particle.
The pressure gradient enforces the volume-preserving constraint and transmits
normal stress. The Laplacian \(\nu\Delta u\) diffuses momentum. The body force
adds external acceleration. This interpretation prevents a common mistake:
pressure is not specified independently in incompressible flow; it is solved
simultaneously with \(u\) so that \(\nabla\cdot u=0\) remains true.
Taking the divergence of the incompressible momentum equation makes this
constraint role explicit:
\[
  -\Delta p
  =
  \rho\,\nabla\cdot(u\cdot\nabla u)-\rho\,\nabla\cdot b,
\]
with boundary conditions determined by the normal component of the momentum
equation at the wall. The pressure is therefore a nonlocal field: it adjusts
instantaneously within the incompressible model, the formal \(c_s\to\infty\)
limit, to project the acceleration back onto divergence-free velocity fields.
Real fluids transmit pressure disturbances at the speed of sound;
incompressibility is the low-Mach, long-wavelength approximation in which that
propagation time is neglected.

\section{Initial And Boundary Data For Navier-Stokes}
\label{sec:navier-stokes-data}

A Navier-Stokes problem is not specified by the differential equation alone.
One must also state the initial velocity, the domain, and the boundary
conditions. For incompressible flow in a fixed domain \(\Omega\), a typical
initial condition is
\[
  u(x,0)=u_0(x),
  \qquad
  \nabla\cdot u_0=0.
\]
The divergence-free condition is not optional: it is the compatibility
condition required by incompressibility. Boundary data must also be compatible
with mass conservation. If \(\partial\Omega\) is an impermeable wall moving
with velocity \(U_{\mathrm{wall}}\), then the normal velocity must satisfy
\[
  u\cdot n=U_{\mathrm{wall}}\cdot n.
\]

The most common viscous wall condition is no slip:
\[
  \boxed{u=U_{\mathrm{wall}}\quad\text{on }\partial\Omega.}
\]
It says that the fluid velocity equals the velocity of the solid boundary, both
normal and tangential components included. No slip is a molecular-scale
constitutive fact about many ordinary fluids at ordinary solid walls; it is not
implied by incompressibility alone. By contrast, inviscid Euler flow uses only
no penetration,
\[
  u\cdot n=U_{\mathrm{wall}}\cdot n,
\]
because tangential slip is not controlled by a viscous boundary layer in the
ideal equation.

Other boundary conditions model different physical surfaces. A free-slip or
stress-free boundary has
\[
  u\cdot n=0,
  \qquad
  (I-n\otimes n)\sigma n=0,
\]
meaning there is no normal flow and no tangential viscous traction. Periodic
boundary conditions identify opposite faces of a box and are useful for
studying homogeneous turbulence and exact periodic solutions such as the
Taylor-Green vortex below.

Pressure boundary conditions deserve special care. For incompressible flow,
pressure is not prescribed independently at a no-slip wall. Instead, taking the
normal component of
\[
  \rho(\partial_tu+u\cdot\nabla u)
  =
  -\nabla p+\mu\Delta u+\rho b
\]
gives the wall relation
\[
  \partial_n p
  =
  n\cdot
  \left(
    \mu\Delta u+\rho b-\rho(\partial_tu+u\cdot\nabla u)
  \right).
\]
Together with the pressure Poisson equation, this determines pressure up to an
additive constant in a closed incompressible domain. This is why numerical
projection methods first advance a tentative velocity and then solve a Poisson
problem for the pressure correction that restores \(\nabla\cdot u=0\).

\section{Nondimensionalization And Reynolds Number}
\label{sec:reynolds-number}

Let \(U\) be a typical speed, \(L\) a typical length, and define
\[
  x=Lx',\qquad
  t=\frac{L}{U}t',\qquad
  u=Uu',\qquad
  p=\rho U^2p'.
\]
Dropping primes, the incompressible Navier-Stokes equation without body force
becomes
\[
  \partial_tu+u\cdot\nabla u
  =
  -\nabla p+\frac1{\mathrm{Re}}\Delta u,
  \qquad
  \nabla\cdot u=0,
\]
where
\[
  \boxed{\mathrm{Re}=\frac{UL}{\nu}},\qquad \nu=\frac{\mu}{\rho}.
\]
The Reynolds number compares inertial transport to viscous diffusion. Small
\(\mathrm{Re}\) flow is viscously dominated and often reversible in time after
the velocity is reversed. Large \(\mathrm{Re}\) flow can carry disturbances
long distances before viscosity damps them.

\section{Energy Identity And Viscous Dissipation}
\label{sec:fluid-energy-identity}

For a smooth incompressible solution in a fixed domain \(\Omega\), assume either
periodic boundary conditions or no-slip walls \(u=0\) on \(\partial\Omega\).
Dot the Navier-Stokes equation with \(u\) and integrate:
\[
  \rho\int_\Omega u\cdot\partial_tu\,\dd^3x
  +
  \rho\int_\Omega u\cdot(u\cdot\nabla u)\,\dd^3x
  =
  -\int_\Omega u\cdot\nabla p\,\dd^3x
  +
  \mu\int_\Omega u\cdot\Delta u\,\dd^3x.
\]
The nonlinear term is a divergence:
\[
  u\cdot(u\cdot\nabla u)=u\cdot\nabla\left(\frac12\norm u^2\right)
  =
  \nabla\cdot\left(\frac12\norm u^2u\right),
\]
because \(\nabla\cdot u=0\). The pressure term is also a divergence:
\[
  u\cdot\nabla p=\nabla\cdot(pu).
\]
The boundary conditions remove both boundary fluxes. Integration by parts gives
\[
  \int_\Omega u\cdot\Delta u\,\dd^3x
  =
  -\int_\Omega |\nabla u|^2\,\dd^3x.
\]
Hence
\[
  \boxed{
  \frac{\dd}{\dd t}
  \int_\Omega \frac12\rho\norm u^2\,\dd^3x
  =
  -\mu\int_\Omega |\nabla u|^2\,\dd^3x\leq0.
  }
\]
Viscosity irreversibly converts organized kinetic energy into heat. This
identity is also the first stability estimate for laminar flows: disturbances
must draw energy from mean shear fast enough to beat viscous dissipation.

\section{Vorticity And Circulation}
\label{sec:vorticity-circulation}

Helmholtz's 1858 vortex theorems and Kelvin's 1869 circulation theorem are the
classical anchors for this section. They explain which parts of inviscid
vorticity motion survive without viscosity.

The vorticity is
\[
  \omega=\nabla\times u.
\]
Taking the curl of the incompressible Navier-Stokes equation gives
\[
  \boxed{
  \partial_t\omega+u\cdot\nabla\omega
  =
  \omega\cdot\nabla u+\nu\Delta\omega.
  }
\]
The term \(\omega\cdot\nabla u\) stretches and tilts vortex lines in three
dimensions. In two dimensions, vorticity is perpendicular to the plane and the
stretching term vanishes:
\[
  \partial_t\omega+u\cdot\nabla\omega=\nu\Delta\omega.
\]
This contrast is central. Two-dimensional inviscid vorticity is transported
like a scalar dye, which creates many conserved quantities. Three-dimensional
vorticity can be stretched by the flow itself, allowing vortex intensification
and making the analysis of Navier-Stokes much harder.
For \(\nu=0\), vorticity is transported by the flow.

\begin{proposition}[Kelvin circulation theorem]
Let \(\gamma(t)\) be a closed material loop moving with an inviscid barotropic
fluid under conservative body forces. Then
\[
  \frac{\dd}{\dd t}\oint_{\gamma(t)}u\cdot \dd x=0.
\]
\end{proposition}

\begin{proof}
Parametrize the loop by \(s\mapsto x(s,t)\), with
\[
  \partial_t x(s,t)=u(x(s,t),t).
\]
Then
\[
  \frac{\dd}{\dd t}\oint u\cdot \partial_sx\,\dd s
  =
  \oint
  \left(
    \frac{D u}{D t}\cdot \partial_sx
    +
    u\cdot\partial_su
  \right)\dd s.
\]
For inviscid barotropic flow with conservative potential \(\Phi\),
\[
  \frac{D u}{D t}=-\nabla h-\nabla\Phi,
\]
where \(h\) is specific enthalpy. The integrand becomes
\[
  \partial_s\left(
    -h-\Phi+\frac12\norm u^2
  \right).
\]
The integral of a single-valued total derivative around a closed loop is zero.
\end{proof}

\begin{proposition}[Helmholtz vortex theorems, inviscid barotropic form]
For a smooth inviscid barotropic fluid with conservative body force:
\begin{enumerate}
\item vortex lines are transported by the fluid motion;
\item the circulation around a small material loop moving with a vortex tube is
constant, so the vortex-tube strength is constant;
\item vortex lines cannot begin or end in the interior of the fluid.
\end{enumerate}
\end{proposition}

\begin{proof}
Kelvin's theorem gives conservation of circulation around every material loop.
Apply it to a small loop encircling a vortex tube. By Stokes' theorem,
\[
  \oint_{\partial S(t)}u\cdot\dd x
  =
  \int_{S(t)}\omega\cdot n\,\dd S,
\]
so the vorticity flux through a material cross-section of the tube is constant.
If a material curve is initially tangent to \(\omega\), then the inviscid
vorticity equation
\[
  \frac{D\omega}{Dt}=\omega\cdot\nabla u-\omega\,\nabla\cdot u
\]
has the same stretching form as the evolution equation for an infinitesimal
material line element per unit mass; hence vortex lines move with the fluid.
Finally, \(\nabla\cdot\omega=\nabla\cdot(\nabla\times u)=0\), so vortex flux has
no interior sources or sinks. Vortex lines either close, meet a boundary, or
continue to infinity.
\end{proof}

\section{Bernoulli's Theorem As An Inviscid First Integral}
\label{sec:bernoulli-theorem}

Daniel Bernoulli's \emph{Hydrodynamica} (1738) is older than the modern
Navier-Stokes equation, but its central invariant appears cleanly from Euler's
equation. For inviscid barotropic flow with conservative body force
\(-\nabla\Phi\),
\[
  \frac{D u}{D t}=-\nabla h-\nabla\Phi,
\]
where \(h\) is specific enthalpy, defined by \(\dd h=\rho^{-1}\dd p\). Assume
the flow is steady. The vector identity
\[
  (u\cdot\nabla)u
  =
  \nabla\left(\frac12|u|^2\right)-u\times\omega,
  \qquad
  \omega=\nabla\times u,
\]
gives
\[
  \nabla\left(\frac12|u|^2+h+\Phi\right)=u\times\omega.
\]
Dotting with \(u\) removes the right-hand side:
\[
  u\cdot\nabla\left(\frac12|u|^2+h+\Phi\right)=0.
\]
Thus
\[
  \boxed{\frac12|u|^2+h+\Phi=\text{constant along each streamline}.}
\]
If the steady flow is also irrotational on a connected region, then
\(u\times\omega=0\), so the same Bernoulli function is constant throughout that
region. The distinction matters: streamline-wise Bernoulli and global Bernoulli
are different statements with different hypotheses.

\section{How Exact Navier-Stokes Solutions Are Found}
\label{sec:exact-ns-solution-strategy}

The full Navier-Stokes equation is nonlinear because of \(u\cdot\nabla u\).
Exact solutions are possible when symmetry, dimensional reduction, or a special
pressure field makes that nonlinearity vanish or become a pure gradient. The
examples below illustrate three reusable mechanisms.

\begin{enumerate}
\item \emph{Unidirectional shear flow.} If
\[
  u=(V(y,t),0,0),
\]
then \(\nabla\cdot u=0\) and
\[
  u\cdot\nabla u=V(y,t)\partial_x(V(y,t),0,0)=0.
\]
The nonlinear term vanishes because the velocity points in the \(x\)-direction
but varies only in the transverse \(y\)-direction.

\item \emph{Fully developed pipe or channel flow.} If the velocity is axial and
depends only on the transverse coordinates, the same cancellation occurs. The
Navier-Stokes equation reduces to an ordinary differential equation in the
cross-section for steady flow, or to a heat equation for unsteady flow.

\item \emph{Nonlinearity absorbed into pressure.} In some periodic flows,
\((u\cdot\nabla)u\) is the gradient of a scalar. Since incompressible pressure
is a Lagrange multiplier, that gradient can be absorbed into \(p\). The
Taylor-Green vortex is the standard example.
\end{enumerate}

For a unidirectional channel flow \(u=(V(y,t),0,0)\) with pressure
\(p=p(x,t)\), incompressible Navier-Stokes becomes
\[
  \rho\partial_t V
  =
  -\partial_xp+\mu\partial_y^2V.
\]
This single equation is the parent of Couette flow, Poiseuille flow, Stokes'
first problem, and the oscillatory Stokes layer. The boundary conditions decide
which solution it represents.

\section{Laminar Flow Solutions}
\label{sec:laminar-solutions}

Laminar solutions are exact, organized solutions of Navier-Stokes. Stability
theory asks which of these organized solutions persist under disturbance.
Each solution follows from three choices: geometry, boundary conditions, and
pressure gradient or wall motion. Changing any one of these choices changes the
profile and the instability problem.

\subsection{The Couette-Poiseuille Family}
\label{subsec:couette-poiseuille-family}

Consider two infinite plates at \(y=0\) and \(y=h\). The lower plate moves with
speed \(U_0\) in the \(x\)-direction, the upper plate moves with speed \(U_h\),
and a constant pressure gradient
\[
  -\frac{\dd p}{\dd x}=G
\]
is imposed. We seek a steady parallel solution
\[
  u=(V(y),0,0),
  \qquad
  p(x)=p_0-Gx.
\]
The ansatz satisfies incompressibility. It also kills the nonlinear term:
\[
  (u\cdot\nabla)u=V(y)\partial_x(V(y),0,0)=0.
\]
The \(y\)- and \(z\)-momentum equations say that \(p\) is independent of
\(y,z\). The \(x\)-momentum equation is
\[
  0=-\frac{\dd p}{\dd x}+\mu V''(y)
  =
  G+\mu V''(y).
\]
Thus
\[
  V''(y)=-\frac{G}{\mu}.
\]
Integrating twice gives
\[
  V(y)=-\frac{G}{2\mu}y^2+C_1y+C_0.
\]
No slip gives
\[
  V(0)=U_0,\qquad V(h)=U_h.
\]
Solving for \(C_0,C_1\) yields the whole Couette-Poiseuille family:
\[
  \boxed{
  V(y)
  =
  U_0+\frac{U_h-U_0}{h}y
  +
  \frac{G}{2\mu}y(h-y).
  }
\]
The first two terms are wall-driven shear. The last term is pressure-driven
flow. The volume flux per unit span is
\[
  Q
  =
  \int_0^h V(y)\,\dd y
  =
  \frac{h}{2}(U_0+U_h)+\frac{Gh^3}{12\mu}.
  \]
The wall shear stresses are
\[
  \tau(0)=\mu V'(0)
  =
  \mu\frac{U_h-U_0}{h}+\frac{Gh}{2},
  \qquad
  \tau(h)=\mu V'(h)
  =
  \mu\frac{U_h-U_0}{h}-\frac{Gh}{2}.
\]
These formulas are often the easiest way to check signs: pressure forcing
increases shear at one wall and decreases it at the other.

\subsection{Plane Couette Flow}
\label{subsec:plane-couette}

Two infinite plates are separated by distance \(h\). The lower plate is fixed;
the upper plate moves with velocity \(Ue_x\). Seek a steady parallel flow
\[
  u=(V(y),0,0),\qquad p=\text{constant}.
\]
The incompressibility condition is automatic. The \(x\)-momentum equation is
\[
  0=\mu V''(y),
\]
with boundary conditions \(V(0)=0\), \(V(h)=U\). Thus
\[
  \boxed{V(y)=U\frac{y}{h}.}
\]
This is the special case of the Couette-Poiseuille family with
\[
  U_0=0,\qquad U_h=U,\qquad G=0.
\]
The shear stress is constant:
\[
  \sigma_{xy}=\mu V'(y)=\mu\frac{U}{h}.
\]

\subsection{Plane Poiseuille Flow}
\label{subsec:plane-poiseuille}

Now both plates are fixed and a pressure gradient drives the flow. Let
\[
  u=(V(y),0,0),
  \qquad
  -\frac{\dd p}{\dd x}=G>0,
  \qquad
  0<y<h.
\]
The steady \(x\)-momentum equation is
\[
  0=-\frac{\dd p}{\dd x}+\mu V''(y)=G+\mu V''(y),
\]
so
\[
  V''(y)=-\frac{G}{\mu}.
\]
With no-slip conditions \(V(0)=V(h)=0\),
\[
  \boxed{
  V(y)=\frac{G}{2\mu}y(h-y).
  }
\]
This is the special case \(U_0=U_h=0\) of the Couette-Poiseuille family.
The maximum speed occurs at \(y=h/2\):
\[
  V_{\max}=\frac{Gh^2}{8\mu}.
\]
The volume flux per unit span is
\[
  Q=\int_0^h V(y)\,\dd y
  =
  \frac{Gh^3}{12\mu}.
\]

\subsection{Hagen-Poiseuille Pipe Flow}
\label{subsec:hagen-poiseuille}

Hagen (1839) and Poiseuille (1840) measured this pipe-flow law experimentally
before it became a standard derivation from
Navier-Stokes. It is a useful reminder that some continuum laws entered
mechanics through measurement before their modern derivation was settled.

For a circular pipe of radius \(a\), use cylindrical radius \(r\) and seek
\[
  u=V(r)e_z,\qquad -\frac{\dd p}{\dd z}=G>0.
\]
The axial equation is
\[
  0=G+\mu\frac1r\frac{\dd}{\dd r}
  \left(
    r\frac{\dd V}{\dd r}
  \right).
\]
Multiply by \(r\) and integrate:
\[
  \frac{\dd}{\dd r}\left(rV'(r)\right)
  =
  -\frac{G}{\mu}r,
  \qquad
  rV'(r)=-\frac{G}{2\mu}r^2+C_1.
\]
Regularity at the centerline requires \(V'(0)\) finite, so \(C_1=0\). A second
integration gives
\[
  V(r)=-\frac{G}{4\mu}r^2+C_2.
\]
No slip at \(r=a\) fixes \(C_2=Ga^2/(4\mu)\). Therefore
\[
  \boxed{
  V(r)=\frac{G}{4\mu}(a^2-r^2).
  }
\]
The flux is
\[
  Q=2\pi\int_0^a V(r)r\,\dd r
  =
  \boxed{\frac{\pi a^4G}{8\mu}}.
\]
The fourth power in \(a\) is one of the most important practical consequences
of viscous laminar flow.

\begin{figure}[htbp]
\centering
\begin{tikzpicture}[scale=1.0, >=Latex, font=\small]
  \draw[->] (-0.25,0) -- (-0.25,3.25) node[above] {\(y\)};
  \draw[->] (-0.25,0) -- (1.2,0) node[right] {\(x\)};
  \draw[thick] (0,0) -- (3.0,0);
  \draw[thick] (0,3.0) -- (3.0,3.0);
  \draw[->, thick, black!70] (2.02,3.22) -- (2.80,3.22)
    node[midway, above, fill=white, inner sep=1.5pt] {wall speed \(U\)};
  \foreach \y/\len in {0.45/0.25,0.95/0.55,1.45/0.85,1.95/1.15,2.45/1.45}
    \draw[blue, thick, ->] (0.35,\y) -- ++(\len,0);
  \node[blue, anchor=west, fill=white, inner sep=1.5pt] at (0.25,2.76) {Couette};
  \node at (1.5,-0.35) {\(V(0)=0,\ V(h)=U\)};

  \begin{scope}[xshift=5.2cm]
    \draw[thick] (0,0) -- (3.0,0);
    \draw[thick] (0,3.0) -- (3.0,3.0);
    \draw[->, thick, black!70] (0.28,3.25) -- (1.16,3.25)
      node[midway, above, fill=white, inner sep=1.5pt] {\(-\partial_xp=G\)};
    \foreach \y/\len in {0.45/0.20,0.95/0.85,1.50/1.65,2.05/0.85,2.55/0.20}
      \draw[red!75!black, thick, ->] (0.35,\y) -- ++(\len,0);
    \node[red!75!black, fill=white, inner sep=1.5pt] at (2.36,2.55) {Poiseuille};
    \node at (1.5,-0.35) {\(V(0)=V(h)=0\)};
  \end{scope}
\end{tikzpicture}
\caption{Couette flow is linear because it is driven by wall motion. Poiseuille
flow is parabolic because it balances pressure forcing against viscous
diffusion of momentum. The arrows are spatial velocity vectors, not a graph
with velocity on the vertical axis.}
\label{fig:couette-poiseuille-profiles}
\end{figure}

\section{Unsteady Exact Solutions}
\label{sec:unsteady-exact-ns-solutions}

Steady laminar solutions show how pressure and wall motion balance viscosity.
Unsteady exact solutions answer a different physical question: how fast does a
boundary or vortex imprint its motion on the surrounding fluid? The answer is
not a sound-crossing time in the incompressible model. It is the diffusive
viscous time \(y^2/\nu\), and the associated length scale is
\(\sqrt{\nu t}\). The examples below are therefore more than formulae. They
show how momentum diffuses, how oscillatory forcing produces a penetration
depth, and why exact solutions are useful tests for numerical codes that claim
to solve the full Navier-Stokes equation.

\subsection{Stokes' First Problem}
\label{subsec:stokes-first-problem}

Stokes' 1851 first problem is the canonical similarity solution for viscous
diffusion of momentum from an impulsively moved wall.

Let a viscous incompressible fluid occupy the half-space \(y>0\). For
\(t<0\), both fluid and wall are at rest. At \(t=0\), the wall at \(y=0\) is
instantaneously set into motion with velocity \(Ue_x\). Seek
\[
  u=(V(y,t),0,0),
  \qquad
  p=\text{constant}.
\]
As before, incompressibility is automatic and \(u\cdot\nabla u=0\). The
Navier-Stokes equation reduces to the heat equation
\[
  \partial_t V=\nu\,\partial_y^2V.
\]
The data are
\[
  V(y,0)=0\quad(y>0),
  \qquad
  V(0,t)=U,
  \qquad
  V(y,t)\to0\quad(y\to\infty).
\]
The only length scale that can be made from \(\nu\) and \(t\) is
\(\sqrt{\nu t}\), so introduce
\[
  \eta=\frac{y}{2\sqrt{\nu t}},
  \qquad
  V(y,t)=U F(\eta).
\]
Using
\[
  \partial_t\eta=-\frac{\eta}{2t},
  \qquad
  \partial_y\eta=\frac{1}{2\sqrt{\nu t}},
\]
the heat equation becomes
\[
  -\frac{U\eta}{2t}F'(\eta)
  =
  \nu U\frac{1}{4\nu t}F''(\eta),
  \qquad
  F''+2\eta F'=0.
\]
The boundary conditions are \(F(0)=1\) and \(F(\infty)=0\). Hence
\[
  \boxed{
  V(y,t)=U\,\operatorname{erfc}\left(\frac{y}{2\sqrt{\nu t}}\right).
  }
\]
The thickness of the region significantly affected by the wall grows like
\(\sqrt{\nu t}\). The wall shear stress has magnitude
\[
  |\tau_w(t)|
  =
  \mu\left|\partial_yV(0,t)\right|
  =
  \mu\frac{U}{\sqrt{\pi\nu t}},
\]
which is singular at \(t=0^+\) because the idealized wall velocity jumps
instantaneously.

\subsection{Oscillatory Stokes Layer}
\label{subsec:oscillatory-stokes-layer}

Now let the wall at \(y=0\) oscillate tangentially:
\[
  V(0,t)=U\cos(\Omega_s t),
  \qquad
  V(y,t)\to0\quad(y\to\infty).
\]
Again \(V\) solves \(\partial_tV=\nu V_{yy}\). Use a complex solution and take
the real part at the end:
\[
  V(y,t)=\operatorname{Re}\left\{U e^{i\Omega_s t-\lambda y}\right\}.
\]
Substitution gives
\[
  i\Omega_s=\nu\lambda^2.
\]
The root with positive real part is
\[
  \lambda=\frac{1+i}{\delta_s},
  \qquad
  \delta_s=\sqrt{\frac{2\nu}{\Omega_s}}.
\]
Therefore
\[
  \boxed{
  V(y,t)=U e^{-y/\delta_s}
  \cos\left(\Omega_s t-\frac{y}{\delta_s}\right).
  }
\]
The length \(\delta_s\) is the oscillatory Stokes-layer thickness. Faster
oscillation gives a thinner layer; larger viscosity gives a thicker layer. The
phase lag \(y/\delta_s\) records the time delay caused by viscous diffusion of
momentum away from the moving wall.

\subsection{Taylor-Green Vortex}
\label{subsec:taylor-green-vortex}

The Taylor-Green vortex is an exact periodic solution of the two-dimensional
incompressible Navier-Stokes equation. On a periodic square, set
\[
\begin{aligned}
  u_x(x,y,t)&=A(t)\sin(kx)\cos(ky),\\
  u_y(x,y,t)&=-A(t)\cos(kx)\sin(ky),\\
  A(t)&=U_0e^{-2\nu k^2t}.
\end{aligned}
\]
The divergence vanishes because
\[
  \partial_xu_x+\partial_yu_y
  =
  A k\cos(kx)\cos(ky)
  -
  A k\cos(kx)\cos(ky)=0.
\]
The Laplacian satisfies
\[
  \Delta u=-2k^2u,
\]
so \(\partial_tu=\nu\Delta u\). The remaining nonlinear term is not zero, but
it is a gradient:
\[
  (u\cdot\nabla)u
  =
  \frac{A(t)^2k}{2}
  \left(\sin(2kx),\,\sin(2ky)\right).
\]
It is cancelled by the pressure
\[
  \boxed{
  p(x,y,t)=p_0+\frac{\rho A(t)^2}{4}
  \left(\cos(2kx)+\cos(2ky)\right).
  }
\]
Indeed,
\[
  -\frac1\rho\nabla p
  =
  \frac{A(t)^2k}{2}
  \left(\sin(2kx),\,\sin(2ky)\right)
  =
  (u\cdot\nabla)u.
\]
Thus the full nonlinear equation
\[
  \partial_tu+u\cdot\nabla u
  =
  -\frac1\rho\nabla p+\nu\Delta u
\]
is satisfied exactly. The vorticity is
\[
  \omega_z
  =
  \partial_xu_y-\partial_yu_x
  =
  2kA(t)\sin(kx)\sin(ky),
\]
and the cell-averaged kinetic energy density is
\[
  \left\langle \frac12\rho |u|^2\right\rangle
  =
  \frac14\rho A(t)^2
  =
  \frac14\rho U_0^2e^{-4\nu k^2t}.
\]

\begin{figure}[htbp]
\centering
\begin{tikzpicture}[scale=1.0, >=Latex, font=\small]
  \begin{scope}[yscale=2.0]
    \draw[->] (-0.3,0) -- (3.6,0) node[right] {\(y\)};
    \draw[->] (0,-0.12) -- (0,1.22) node[above] {\(V/U\)};
    \draw[thick] (0,0) -- (0,1.0);
    \node[left] at (0,1.0) {moving wall};
    \draw[blue!70!black, thick] plot[smooth] coordinates {
      (0.0,1.000) (0.2,0.778) (0.4,0.572) (0.6,0.396)
      (0.8,0.258) (1.0,0.158) (1.2,0.089) (1.4,0.048)
      (1.6,0.024) (1.8,0.011) (2.0,0.005) (2.2,0.002)
      (2.4,0.001) (2.6,0.000) (2.8,0.000) (3.0,0.000) (3.2,0.000)
    };
    \draw[red!75!black, thick] plot[smooth] coordinates {
      (0.0,1.000) (0.2,0.871) (0.4,0.744) (0.6,0.624)
      (0.8,0.514) (1.0,0.414) (1.2,0.327) (1.4,0.253)
      (1.6,0.192) (1.8,0.142) (2.0,0.102) (2.2,0.072)
      (2.4,0.050) (2.6,0.034) (2.8,0.022) (3.0,0.014) (3.2,0.009)
    };
    \draw[purple!75!black, thick] plot[smooth] coordinates {
      (0.0,1.000) (0.2,0.903) (0.4,0.808) (0.6,0.715)
      (0.8,0.627) (1.0,0.543) (1.2,0.465) (1.4,0.394)
      (1.6,0.330) (1.8,0.274) (2.0,0.224) (2.2,0.180)
      (2.4,0.144) (2.6,0.114) (2.8,0.089) (3.0,0.068) (3.2,0.052)
    };
    \node[blue!70!black, fill=white, inner sep=1pt] at (0.98,0.27) {\(t_1\)};
    \node[red!75!black, fill=white, inner sep=1pt] at (1.70,0.36) {\(t_2\)};
    \node[purple!75!black, fill=white, inner sep=1pt] at (2.45,0.31) {\(t_3\)};
  \end{scope}
  \node[align=center] at (1.8,-0.55) {Stokes first problem:\\viscous penetration grows like \(\sqrt{\nu t}\)};

  \begin{scope}[xshift=6.0cm]
    \draw[->] (-0.2,0) -- (3.2,0) node[right] {\(x\)};
    \draw[->] (0,-0.2) -- (0,3.2) node[above] {\(y\)};
    \draw[step=1.0, gray!35, thin] (0,0) grid (3,3);
    \foreach \x/\y/\dx/\dy in {
      0.5/0.5/0.55/-0.55, 1.5/0.5/0.55/0.55, 2.5/0.5/-0.55/0.55,
      0.5/1.5/-0.55/-0.55, 1.5/1.5/-0.55/0.55, 2.5/1.5/0.55/0.55,
      0.5/2.5/-0.55/0.55, 1.5/2.5/-0.55/-0.55, 2.5/2.5/0.55/-0.55}
      \draw[teal!75!black, thick, ->] (\x,\y) -- ++(\dx,\dy);
    \node[align=center] at (1.5,-0.55) {Taylor-Green vortex:\\periodic eddies decay as \(e^{-2\nu k^2t}\)};
  \end{scope}
\end{tikzpicture}
\caption{Two unsteady exact solutions. Left: an impulsively moved wall creates
the profile \(V/U=\operatorname{erfc}(y/(2\sqrt{\nu t}))\), sampled at
three times. Right: the arrows sample the Taylor-Green velocity
\((A\sin kx\cos ky,-A\cos kx\sin ky)\); advection is absorbed by pressure
while viscosity damps the amplitude.}
\label{fig:unsteady-ns-exact-solutions}
\end{figure}

\section{Linear Stability Of Laminar Shear Flow}
\label{sec:linear-stability-shear}

Exact laminar solutions become physically informative when their nearby motions
are understood. Stability analysis supplies this local question. We linearize
Navier-Stokes around a known base flow and ask whether infinitesimal
disturbances grow or decay. This gives a mathematically controlled diagnostic
that can be compared with experiments and simulations.

Let a steady base flow be
\[
  U(y)e_x.
\]
Write a small perturbation as
\[
  u=U(y)e_x+\tilde u,\qquad
  p=P(x)+\tilde p,
\]
and retain only terms linear in \(\tilde u\). In nondimensional variables,
\[
  \boxed{
  \partial_t\tilde u
  +
  U\partial_x\tilde u
  +
  \tilde v\,U'(y)e_x
  =
  -\nabla\tilde p
  +
  \frac1{\mathrm{Re}}\Delta\tilde u,
  \qquad
  \nabla\cdot\tilde u=0.
  }
\]
The term \(\tilde v\,U'e_x\) is the mechanism by which cross-stream velocity
tilts fluid into faster or slower parts of the mean shear and extracts energy
from the base flow.

\begin{figure}[htbp]
\centering
\begin{tikzpicture}[scale=1.0, >=Latex, font=\small]
  \draw[->] (-0.25,0) -- (-0.25,3.4) node[above] {\(y\)};
  \draw[->] (-0.25,0) -- (5.9,0) node[right] {\(x\)};
  \draw[thick] (0,0.35) -- (5.3,0.35);
  \draw[thick] (0,3.05) -- (5.3,3.05);
  \foreach \y/\len in {0.65/0.35,1.10/0.62,1.55/0.92,2.00/1.25,2.45/1.55,2.85/1.8}
    \draw[blue!70!black, thick, ->] (0.45,\y) -- ++(\len,0);
  \draw[red!75!black, thick, ->] (2.75,1.15) -- (2.75,2.05);
  \node[red!75!black, fill=white, inner sep=1pt] at (2.42,2.16) {\(\tilde v\)};
  \draw[red!75!black, thick, ->] (2.75,2.05) -- (3.55,2.05);
  \draw[purple!75!black, thick, domain=0.3:5.0, samples=120]
    plot (\x,{1.75+0.22*sin(115*\x)});
  \draw[red!75!black, <->] (4.20,1.75) -- (4.20,1.97);
  \node[red!75!black, anchor=west, fill=white, inner sep=1pt] at (4.28,1.86)
    {amp.};
  \node[blue!70!black] at (4.25,2.75) {base shear \(U(y)e_x\)};
  \node[purple!75!black, fill=white, inner sep=1pt] at (4.10,1.16)
    {normal-mode disturbance};
\end{tikzpicture}
\caption{Linear shear-flow stability. A wall-normal perturbation velocity
\(\tilde v\) moves fluid across the base shear; the term
\(\tilde v\,U'(y)e_x\) is the linear energy-exchange channel between the base
flow and the disturbance.}
\label{fig:shear-stability-mechanism}
\end{figure}

\subsection{Orr-Sommerfeld Equation}
\label{subsec:orr-sommerfeld}

Orr and Sommerfeld derived this normal-mode equation independently in
1907--1908. Its history is a useful warning: the equation is linear, but the
transition problem it probes is not exhausted by linear eigenvalues.

For two-dimensional perturbations, introduce a streamfunction
\[
  \tilde u=\partial_y\psi,\qquad
  \tilde v=-\partial_x\psi.
\]
Squire's theorem justifies this restriction for parallel shear flows: if a
three-dimensional normal mode is unstable at Reynolds number \(\mathrm{Re}\),
then there is a two-dimensional mode unstable at a no larger effective Reynolds
number. Thus the least stable linear mode can be found in the two-dimensional
Orr-Sommerfeld problem.
Seek normal modes
\[
  \psi(x,y,t)=\phi(y)e^{i\alpha x+\sigma t}.
\]
Taking the curl of the linearized equation removes pressure and gives the
linear vorticity equation. Since the perturbation vorticity is
\[
  \tilde\omega_z
  =
  \partial_x\tilde v-\partial_y\tilde u
  =
  -(\partial_y^2-\alpha^2)\phi\,e^{i\alpha x+\sigma t},
\]
substitution gives a fourth-order eigenvalue equation for the wall-normal
shape \(\phi(y)\):
\[
  \boxed{
  (\sigma+i\alpha U)(D^2-\alpha^2)\phi
  -
  i\alpha U''\phi
  =
  \frac1{\mathrm{Re}}(D^2-\alpha^2)^2\phi,
  }
\]
where \(D=\dd/\dd y\). No-slip and no-penetration at walls give
\[
  \phi=0,\qquad D\phi=0
\]
at the walls. This fourth-order eigenvalue problem is the
Orr-Sommerfeld equation. A base flow is linearly unstable if some eigenvalue has
\[
  \operatorname{Re}\sigma>0.
\]
For viscous flow on a bounded channel, the no-slip boundary conditions make the
Orr-Sommerfeld operator a fourth-order non-self-adjoint boundary-value problem
with a discrete spectrum. The eigenfunctions are not orthogonal in the way
self-adjoint Sturm-Liouville modes are, and the inviscid Rayleigh limit is
singular: continuous-spectrum effects and critical layers can appear as
\(\mathrm{Re}\to\infty\).
Many fluid-mechanics texts write the same normal mode as
\(e^{i\alpha(x-ct)}\). With the convention used here, the complex phase speed
\(c\) and growth rate \(\sigma\) are related by
\[
  \sigma=-i\alpha c,
\]
so \(\operatorname{Re}\sigma=\alpha\,\operatorname{Im}c\) for real
\(\alpha\).

\begin{remark}[Classical facts]
Plane Poiseuille flow has a linear instability above a critical Reynolds number
about \(5772\), using centerline speed and half-channel width; the computation
was made reliable in the twentieth century by Thomas and later Orszag. Plane
Couette flow is linearly stable for all Reynolds numbers in the classical
Orr-Sommerfeld analysis, yet it can transition through finite-amplitude
disturbances. Reynolds' 1883 dye experiment in pipe flow is the founding
observation behind this warning. Stability is therefore not a single concept:
linear eigenvalue stability, energy stability, transient growth, and nonlinear
transition answer different questions.
\end{remark}

\subsection{Energy Stability Estimate}
\label{subsec:energy-stability}

For perturbations of a parallel flow, the kinetic energy of the disturbance
satisfies
\[
  \frac12\frac{\dd}{\dd t}\int_\Omega |\tilde u|^2\,\dd^3x
  =
  -\int_\Omega \tilde u\,\tilde v\,U'(y)\,\dd^3x
  -
  \frac1{\mathrm{Re}}
  \int_\Omega |\nabla\tilde u|^2\,\dd^3x.
\]
The second term is viscous dissipation. The first is production by mean shear.
If dissipation dominates production for all admissible perturbations, the base
flow is energy stable. If not, the estimate does not prove instability, but it
reveals the only possible energy source for perturbation growth.

\section{Point Vortices As A Hamiltonian Fluid Model}
\label{sec:point-vortices}

In two dimensions, an idealized vorticity distribution can be written as
\[
  \omega(x,y)=\sum_{i=1}^{N}\Gamma_i\delta(x-x_i,y-y_i),
\]
where \(\Gamma_i\) is circulation. Helmholtz introduced point-vortex
idealizations in his 1858 vortex work, and Kirchhoff gave the Hamiltonian form.
Here they serve as a finite-dimensional Hamiltonian reduction of ideal-fluid
motion. The point vortex equations are
\[
\begin{aligned}
  \Gamma_i\dot x_i &= \frac{\partial H}{\partial y_i},\\
  \Gamma_i\dot y_i &= -\frac{\partial H}{\partial x_i},
\end{aligned}
\]
with Hamiltonian
\[
  H
  =
  -\frac{1}{4\pi}
  \sum_{i<j}\Gamma_i\Gamma_j
  \log\left((x_i-x_j)^2+(y_i-y_j)^2\right).
\]
This is Hamiltonian mechanics, but the symplectic form is weighted:
\[
  \omega_{\mathrm{vortex}}
  =
  \sum_i \Gamma_i\,\dd x_i\wedge\dd y_i.
\]
When the circulations have mixed signs this form is not positive in any metric
sense. That is not a problem: symplectic structure requires nondegeneracy, not
positive definiteness.
The weight \(\Gamma_i\) is the circulation carried by the \(i\)-th vortex. One
can derive this finite-dimensional symplectic form by reducing the ideal-fluid
Hamiltonian bracket to delta-function vorticity; in this section we use the
resulting point-vortex bracket as the model.

\begin{figure}[htbp]
\centering
\begin{tikzpicture}[scale=1.05, >=Latex, font=\small]
  \draw[->] (-2.5,0) -- (2.7,0) node[right] {\(x\)};
  \draw[->] (0,-2.2) -- (0,2.4) node[above] {\(y\)};
  \draw[thick] (-0.65,-0.35) circle (0.16);
  \node at (-0.65,-0.35) {\(+\)};
  \node[fill=white, inner sep=1pt] (plusgamma) at (-1.24,-0.76) {\(+\Gamma\)};
  \draw[thin] (plusgamma.north east) -- (-0.77,-0.43);
  \draw[thick] (0.65,0.35) circle (0.16);
  \node at (0.65,0.35) {\(-\)};
  \node[fill=white, inner sep=1pt] (minusgamma) at (1.58,0.16) {\(-\Gamma\)};
  \draw[thin] (minusgamma.west) -- (0.77,0.43);
  \draw[dashed] (-0.65,-0.35) -- (0.65,0.35);
  \draw[<->] (-0.65,-0.70) -- (0.65,0.00)
    node[midway, below right, fill=white, inner sep=1pt] {fixed separation \(d\)};
  \draw[blue, thick, ->] (-2.02,-1.22) -- (-2.52,-0.28);
  \draw[blue, thick, ->] (2.12,-0.10) -- (1.60,0.86);
  \node[blue, align=center, fill=white, inner sep=1pt] at (1.72,1.20)
    {common\\translation};
\end{tikzpicture}
\caption{A vortex dipole is the simplest point-vortex motion: the two vortices
translate together with fixed separation. More complicated point-vortex
systems give finite-dimensional Hamiltonian models of ideal-fluid motion.}
\label{fig:vortex-dipole}
\end{figure}

\section{What To Take Forward}
\label{sec:fluid-take-forward}

\begin{itemize}
\item The material derivative is the bridge between following particles and
describing fields at fixed spatial points.
\item Conservation of mass and momentum are balance laws before they are
differential equations; the local PDEs come from applying the divergence theorem
to arbitrary moving volumes.
\item Navier-Stokes is Cauchy momentum balance plus the Newtonian stress law.
Changing the constitutive law changes the continuum model.
\item Exact Navier-Stokes solutions become tractable when symmetry removes the
nonlinear term, when the flow is fully developed, or when advection can be
absorbed into pressure.
\item The Reynolds number compares inertial transport with viscous diffusion of
momentum. It is the organizing scale behind laminar solutions, instabilities,
and transition.
\item Point vortices show that Hamiltonian structure survives inside fluid
mechanics after a strong but explicit reduction.
\end{itemize}

\paragraph{Questions to carry forward.}
Which part of Navier-Stokes is kinematics, which is balance law, and which is
constitutive assumption? When is pressure a thermodynamic variable, and when is
it a constraint multiplier? Which laminar solutions remain stable under small
disturbances?

\section{Chapter-End Laboratories And Exercises}
\label{sec:fluid-chapter-end-labs}

The following derivations isolate two recurring fluid-mechanics diagnostics:
pressure as the multiplier that enforces incompressibility, and the
Orr-Sommerfeld eigenvalue problem as the linear stability test for parallel
shear flows.

\subsection{Pressure As A Lagrange Multiplier In Incompressible Flow}
\label{sec:detailed-pressure-multiplier}

For incompressible Navier-Stokes,
\[
  \partial_tu+u\cdot\nabla u
  =
  -\frac1\rho\nabla p+\nu\Delta u,
  \qquad
  \nabla\cdot u=0.
\]
The pressure is not determined by a separate thermodynamic equation in the
incompressible model. It is the Lagrange multiplier enforcing
\(\nabla\cdot u=0\).

Take the divergence of the momentum equation. Since
\[
  \nabla\cdot\partial_tu=\partial_t(\nabla\cdot u)=0
  \quad\text{and}\quad
  \nabla\cdot\Delta u=\Delta(\nabla\cdot u)=0,
\]
this gives the pressure Poisson equation
\[
  \boxed{
  -\frac1\rho\Delta p
  =
  \nabla\cdot(u\cdot\nabla u)
  =
  \partial_i u_j\,\partial_j u_i.
  }
\]
The last equality follows by expanding:
\[
  \partial_i(u_j\partial_j u_i)
  =
  (\partial_i u_j)(\partial_j u_i)
  +
  u_j\partial_j(\partial_i u_i),
\]
and the second term vanishes by incompressibility. Thus \(p\) is whatever
spatial field is needed to project the nonlinear acceleration back onto the
space of divergence-free vector fields.

This is the Helmholtz projection in PDE form. On a smooth bounded or periodic
domain, with the appropriate boundary conditions, a vector field \(a\) can be
decomposed as
\[
  a=a_{\mathrm{df}}+\nabla\phi,\qquad \nabla\cdot a_{\mathrm{df}}=0.
\]
Taking the divergence gives a Poisson equation
\[
  \Delta\phi=\nabla\cdot a,
\]
with boundary conditions determined by the normal component. In incompressible
Navier-Stokes, the pressure supplies exactly this gradient part of the
acceleration.

\subsection{Orr-Sommerfeld Boundary Conditions And Spectrum}
\label{sec:detailed-orr-sommerfeld}

The body of the chapter derived the Orr-Sommerfeld equation. The laboratory
here records the boundary and spectral data needed to turn that differential
equation into a well-posed eigenvalue problem. No penetration means
\[
  \tilde v=-\partial_x\psi=0
  \quad\Rightarrow\quad
  \phi=0
\]
at a wall. No slip also requires
\[
  \tilde u=\partial_y\psi=0
  \quad\Rightarrow\quad
  D\phi=0.
\]
Thus the Orr-Sommerfeld equation is fourth order and needs the four wall
conditions \(\phi=D\phi=0\) at the two boundaries.
For a channel \(y\in[-1,1]\), the eigenvalue problem may be written as
\[
  \mathcal L_{\mathrm{OS}}\phi=\sigma\,\mathcal M\phi,
  \qquad
  \mathcal M=D^2-\alpha^2,
\]
where
\[
  \mathcal L_{\mathrm{OS}}
  =
  -i\alpha U(D^2-\alpha^2)+i\alpha U''
  +\frac1{\mathrm{Re}}(D^2-\alpha^2)^2.
\]
The no-slip boundary conditions make this a discrete non-self-adjoint spectral
problem for finite \(\mathrm{Re}\). Numerically, one must monitor convergence
with respect to the wall-normal discretization because non-normal operators can
produce spurious eigenvalues and large transient amplification even when every
eigenvalue is stable.

\section{Associated Demos}
\label{sec:fluid-associated-demos}

\begin{itemize}
\item \path{demos/python/navier_stokes_solutions.py}.
Exact Couette-Poiseuille, pipe, Stokes-layer, and Taylor-Green solutions.
\item \path{demos/python/fluids_vorticity.py}, a point-vortex Hamiltonian model.
\item \path{demos/python/hamiltonian_pendulum.py}, useful for comparing
separatrix geometry with shear-flow stability.
\item Suggested extension: a numerical Orr-Sommerfeld spectrum for Couette and
Poiseuille flow.
\end{itemize}
